\section{Actual proper-center line frames in the original completed point field}

\begin{definition}\label{def:fc-point-completed-line-frame}
\lean{EnvelopeProperty.Berkovich.SchemePointCompletedLineFrame.completedFrame,EnvelopeProperty.Berkovich.SchemePointCompletedLineFrame.completedLocalFrame}
\leanok
\uses{cor:actual-line-model-norm-local-frame,scheme:actual-native-completed-residue-field,def:fc-point-laurent-line-evaluation}
Let $p:X\to\operatorname{Spec}R$ be proper, let $x$ be an original
relative valued point, and suppose the base valuation is bounded by
one. Let $U$ be an original open neighborhood of its actual proper
center, with an actual frame of the original line $L$ on $U$.
The actual residue-field frame constructed from this neighborhood
pulls back through the original residue-field embedding into $H(x)$
to an actual frame of the same original completed-point pullback
$h^*L$. Use this frame in the preceding original linear functional.
\end{definition}
\begin{proof}
The original residue-field morphism factors through $U$ by the actual
proper-center construction. Pull the given original frame along
this morphism, obtaining the entire original residue-field frame.
The original completed field point is the composite of this residue
morphism with the spectrum of the actual residue-field embedding.
The native composition and structure-unit comparisons therefore
transport the residue frame to the stated completed-point frame.
No completion-frame compatibility is assumed.
\end{proof}

\begin{theorem}\label{thm:fc-point-completed-frame-model-norm}
\lean{EnvelopeProperty.Berkovich.SchemePointCompletedLineFrame.completedFrame_functional,EnvelopeProperty.Berkovich.SchemePointCompletedLineFrame.completedLocalFrame_functional,EnvelopeProperty.Berkovich.SchemePointCompletedLineFrame.completedLocalFrame_nnnorm_eq_modelLineNorm}
\leanok
\uses{def:fc-point-completed-line-frame,cor:actual-line-model-norm-local-frame,lem:fc-original-line-section-fibre,scheme:actual-native-completed-residue-field}
For every original global section $m$ of an original invertible line,
the absolute value of its original completed-frame functional is
exactly the true original proper-model line norm of the original
residue-fibre section of $m$. This includes the zero section.
\end{theorem}
\begin{proof}
Apply the actual residue frame to the native pullback-unit section
of $m$, and denote its original scalar by $c\in\kappa_X(x)$.
The proper-center local-frame theorem computes the true model
norm as the original residue valuation of $c$. Pullback of the
whole original scalar section morphism through the residue-field
embedding gives precisely the embedded $c$ as the completed-frame
functional. This follows from the genuine native pullback-unit,
composition and scalar formulas, not from a supplied evaluation
identity. The original residue embedding into $H(x)$ preserves its
absolute value. These equalities prove the assertion for every
original section, including $c=0$.
\end{proof}

\begin{corollary}\label{cor:fc-point-exists-completed-frame-model-norm}
\lean{EnvelopeProperty.Berkovich.SchemePointCompletedLineFrame.exists_completedFrame_modelLineNorm,EnvelopeProperty.Berkovich.SchemePointCompletedLineFrame.exists_completedFrame_finite_family_modelNorm}
\leanok
\uses{thm:fc-point-completed-frame-model-norm,def:fc-point-completed-line-frame,thm:fc-point-laurent-line-gauss}
Every original invertible line has an actual completed-point frame
for which the preceding equality holds simultaneously for all
original global sections. Consequently, at the actual original
Laurent Gauss source point, every actual original Laurent-pulled
line section has the weighted maximum of the true original model
norms of its unique original coefficients. For a finite original
section family, the same original finite coefficient and integer
weight family is chosen before all original valued points and
positive radii.
\end{corollary}
\begin{proof}
The actual open frame cover of the invertible line contains its
actual proper center in one member. Choose that original member
and frame, and apply the preceding theorem to every original
section. Substituting this equality into the genuine native
product-frame Gauss coefficient formula gives the weighted
original model norms. The original coefficient-pair construction
depends only on the original section family; it precedes all
original points, frames and radii. No boundedness of the Laurent
Gauss base valuation is used.
\end{proof}
